% ================================================================
%  mermaid-showcase.tex
%  A richly illustrated document using mermaid-overleaf.sty
%  hosted on GitHub.
%
%  HOW TO COMPILE ON OVERLEAF
%  ─────────────────────────────────────────────────────────────
%  1. Create a new Overleaf project.
%  2. In the left sidebar click  New File  →  From External URL
%     and paste:
%       https://raw.githubusercontent.com/YOUR-USERNAME/mermaid-overleaf/main/mermaid-overleaf.sty
%     Save it as  mermaid-overleaf.sty
%  3. Upload this file (mermaid-showcase.tex) as your main file.
%  4. Set the compiler to  pdfLaTeX  and hit Compile.
%     Everything renders without shell-escape or any extra tools.
%
%  To promote a listing to a real rendered diagram:
%    • Paste the Mermaid source into https://mermaid.live
%    • Export as PDF → upload to Overleaf → add  image={file.pdf}
%      to the environment options.
% ================================================================

\documentclass[12pt, a4paper]{article}

% ── Load mermaid-overleaf from your GitHub repo ──────────────
% (After adding it as a linked file via "From External URL")
\usepackage[theme=ocean]{mermaid-overleaf}

% ── Document packages ─────────────────────────────────────────
\usepackage[T1]{fontenc}
\usepackage[utf8]{inputenc}
\usepackage{lmodern}
\usepackage{microtype}
\usepackage[margin=2.2cm, top=2.5cm, bottom=2.8cm]{geometry}
\usepackage{hyperref}
\usepackage{booktabs}
\usepackage{array}
\usepackage{parskip}
\usepackage{titlesec}
\usepackage{fancyhdr}
\usepackage{xcolor}

% ── Hyperref colours ──────────────────────────────────────────
\hypersetup{
  colorlinks = true,
  linkcolor  = {mmd-ocean-frame},
  urlcolor   = {mmd-ocean-accent},
  citecolor  = {mmd-forest-frame},
  pdftitle   = {Mermaid Diagrams Showcase},
  pdfauthor  = {mermaid-overleaf},
}

% ── Section styling ───────────────────────────────────────────
\titleformat{\section}
  {\large\bfseries\sffamily\color{mmd-ocean-frame}}
  {\thesection}{1em}{}[\color{mmd-ocean-accent}\titlerule]

\titleformat{\subsection}
  {\normalsize\bfseries\sffamily\color{mmd-forest-frame}}
  {\thesubsection}{0.8em}{}

% ── Header / footer ───────────────────────────────────────────
\pagestyle{fancy}
\fancyhf{}
\renewcommand{\headrulewidth}{0.4pt}
\fancyhead[L]{\small\sffamily\color{mmd-ocean-frame} Mermaid Diagrams Showcase}
\fancyhead[R]{\small\sffamily\color{mmd-ocean-cmt} \texttt{mermaid-overleaf} · GitHub}
\fancyfoot[C]{\small\thepage}
\renewcommand{\headrule}{\color{mmd-ocean-accent}\hrule width\headwidth height\headrulewidth}

% ── Convenience macro: coloured term in prose ─────────────────
\newcommand{\term}[1]{{\sffamily\bfseries\color{mmd-ocean-kw}#1}}
\newcommand{\pkg}[1]{{\ttfamily\color{mmd-ocean-accent}#1}}

% =============================================================
\begin{document}
% =============================================================

% ── Title page ───────────────────────────────────────────────
\begin{titlepage}
  \centering
  \vspace*{3cm}

  {\fontsize{36}{42}\selectfont\sffamily\bfseries
   \textcolor{mmd-ocean-frame}{Mermaid}
   \textcolor{mmd-ocean-accent}{Diagrams}\\[0.3em]
   \textcolor{mmd-forest-frame}{Showcase}}

  \vspace{1em}
  {\large\sffamily\color{mmd-ocean-cmt}
    A complete gallery of diagram types, themes, and use-cases\\
    compiled entirely on cloud Overleaf}

  \vspace{2.5cm}

  \begin{tcolorbox}[
      enhanced,
      colback=mmd-ocean-bg,
      colframe=mmd-ocean-frame,
      boxrule=1.5pt,
      arc=8pt,
      width=0.72\textwidth,
      halign=center,
      fontupper=\sffamily,
  ]
    {\color{mmd-ocean-kw}\bfseries Package:}\enspace
    {\color{mmd-ocean-str}\ttfamily mermaid-overleaf}\\[4pt]
    {\color{mmd-ocean-kw}\bfseries Hosted at:}\enspace
    {\color{mmd-ocean-accent}\ttfamily github.com/YOUR-USERNAME/mermaid-overleaf}\\[4pt]
    {\color{mmd-ocean-kw}\bfseries Themes:}\enspace
    {\color{mmd-ocean-str}ocean · forest · sunset · arctic · midnight}\\[4pt]
    {\color{mmd-ocean-kw}\bfseries Compiled with:}\enspace
    {\color{mmd-ocean-str}pdfLaTeX on Overleaf — no shell-escape}
  \end{tcolorbox}

  \vspace{2cm}
  {\small\sffamily\color{mmd-ocean-cmt} \today}

  \vfill

  \MermaidWorkflow
\end{titlepage}

% ── Table of contents ─────────────────────────────────────────
\tableofcontents
\newpage

% =============================================================
\section{Flowcharts}
% =============================================================

Flowcharts describe process logic with nodes and directed edges.
The \pkg{mermaid-overleaf} package renders the source in a
syntax-highlighted listing with a coloured frame — switch the
\term{theme} option per diagram for variety.

% ── 1. CI/CD pipeline ────────────────────────────────────────
\begin{mermaid}[
    type    = flowchart,
    theme   = ocean,
    caption = {Continuous integration and deployment pipeline},
    label   = {fig:cicd},
    width   = 0.92\linewidth
]
flowchart TD
    A([Developer Push]) --> B[GitHub Actions triggered]
    B --> C{Branch?}

    C -->|main| D[Run full test suite]
    C -->|feature/*| E[Run unit tests only]
    C -->|hotfix/*| F[Run smoke tests]

    D --> G{All tests pass?}
    E --> G
    F --> G

    G -->|No| H[/Notify team via Slack/]
    G -->|Yes| I[Build Docker image]

    I --> J[Push to container registry]
    J --> K{Environment}

    K -->|Staging| L[Deploy to staging]
    K -->|Production| M{Manual approval?}

    L --> N[Run E2E tests]
    N -->|Pass| O([Ready for prod])
    N -->|Fail| H

    M -->|Approved| P[Blue-green deploy]
    M -->|Rejected| Q([Hold])
    P --> R[Health check]
    R -->|OK| S([Live in production])
    R -->|Fail| T[Rollback]
    T --> H
\end{mermaid}

% ── 2. Microservice data flow ─────────────────────────────────
\begin{mermaid}[
    type    = flowchart,
    theme   = midnight,
    caption = {E-commerce microservice data-flow — left to right},
    label   = {fig:microservice},
    width   = \linewidth
]
flowchart LR
    subgraph Clients["🌐 Clients"]
        WEB[Web Browser]
        MOB[Mobile App]
        EXT[Third-party API]
    end

    subgraph Edge["⚡ Edge Layer"]
        GW[API Gateway]
        LB[Load Balancer]
        RL[Rate Limiter]
    end

    subgraph Core["🔧 Core Services"]
        AUTH[Auth Service]
        PROD[Product Service]
        ORD[Order Service]
        PAY[Payment Service]
        NOTIF[Notification Service]
    end

    subgraph Data["🗄 Data Layer"]
        PGAUTH[(Auth DB)]
        PGPROD[(Product DB)]
        PGORD[(Order DB)]
        REDIS[(Redis Cache)]
        S3[(S3 Storage)]
        QUEUE[[Message Queue]]
    end

    WEB & MOB & EXT --> GW
    GW --> LB --> RL
    RL --> AUTH & PROD & ORD
    AUTH --> PGAUTH
    PROD --> PGPROD & REDIS & S3
    ORD --> PGORD & REDIS
    ORD --> PAY
    PAY --> QUEUE
    QUEUE --> NOTIF
    ORD --> QUEUE
\end{mermaid}

% ── 3. Decision tree ──────────────────────────────────────────
\begin{mermaid}[
    type    = flowchart,
    theme   = forest,
    caption = {Machine learning model selection decision tree},
    label   = {fig:mlselect},
    width   = 0.88\linewidth
]
flowchart TD
    S([Start: choose a model]) --> Q1{Labelled data available?}

    Q1 -->|Yes| Q2{Problem type?}
    Q1 -->|No|  Q3{Goal?}

    Q2 -->|Classification| Q4{Large dataset?}
    Q2 -->|Regression|     Q5{Non-linear?}
    Q2 -->|Sequence|       RNN[RNN / LSTM / Transformer]

    Q4 -->|Yes| RF[Random Forest\nor XGBoost]
    Q4 -->|No|  SVM[SVM or Logistic Regression]

    Q5 -->|Yes| NN[Neural Network\nRegressor]
    Q5 -->|No|  LR[Linear Regression\nor Ridge]

    Q3 -->|Discover structure| CL{High dimensions?}
    Q3 -->|Generate data|      GEN[GAN / VAE / Diffusion]
    Q3 -->|Learn to act|       RL[Reinforcement Learning]

    CL -->|Yes| AE[Autoencoder\nor t-SNE]
    CL -->|No|  KM[K-Means\nor DBSCAN]
\end{mermaid}

% =============================================================
\section{Sequence Diagrams}
% =============================================================

Sequence diagrams capture time-ordered interactions between
participants. They are especially valuable for documenting APIs,
authentication flows, and distributed-system protocols.

% ── 4. JWT auth flow ─────────────────────────────────────────
\begin{mermaid}[
    type    = sequence,
    theme   = sunset,
    caption = {JWT authentication and token-refresh flow},
    label   = {fig:jwt},
    width   = 0.92\linewidth
]
sequenceDiagram
    actor     U  as 👤 User
    participant C  as 🖥 Client App
    participant A  as 🔐 Auth Server
    participant R  as 🗃 Resource API
    participant DB as 🗄 Database

    U  ->> C  : Enter credentials
    C  ->> A  : POST /auth/login {email, password}
    A  ->> DB : SELECT user WHERE email=?
    DB -->> A : User record + hashed password

    alt Credentials valid
        A  ->> A  : Generate access_token (15 min)
        A  ->> A  : Generate refresh_token (7 days)
        A  -->> C : 200 OK {access_token, refresh_token}
        C  ->> C  : Store tokens securely
        C  -->> U : Dashboard
    else Invalid credentials
        A  -->> C : 401 Unauthorized
        C  -->> U : Show error message
    end

    Note over C,R: Later — making an API call

    C  ->> R  : GET /api/orders (Bearer: access_token)
    R  ->> A  : Verify token signature
    A  -->> R : Token valid, claims: {userId, roles}
    R  ->> DB : SELECT orders WHERE userId=?
    DB -->> R : Order rows
    R  -->> C : 200 OK {orders: [...]}

    Note over C,A: Token expires — silent refresh

    C  ->> A  : POST /auth/refresh {refresh_token}
    A  ->> A  : Verify refresh token, rotate it
    A  -->> C : 200 OK {new access_token, new refresh_token}
\end{mermaid}

% ── 5. Distributed transaction (Saga pattern) ─────────────────
\begin{mermaid}[
    type    = sequence,
    theme   = midnight,
    caption = {Saga pattern — distributed order transaction with compensations},
    label   = {fig:saga},
    width   = \linewidth
]
sequenceDiagram
    participant ORC as 🎭 Saga Orchestrator
    participant ORD as 📦 Order Service
    participant INV as 🏭 Inventory Service
    participant PAY as 💳 Payment Service
    participant SHP as 🚚 Shipping Service

    ORC ->> ORD : createOrder(items, userId)
    ORD -->> ORC : orderId=7842, status=PENDING

    ORC ->> INV : reserveItems(orderId, items)

    alt Items available
        INV -->> ORC : reserved=true, reservationId=R99

        ORC ->> PAY : chargeCustomer(orderId, amount)

        alt Payment succeeds
            PAY -->> ORC : charged=true, txId=TX-44

            ORC ->> SHP : scheduleShipment(orderId)
            SHP -->> ORC : shipmentId=SH-21, eta=2d

            ORC ->> ORD : updateStatus(CONFIRMED)
            ORD -->> ORC : ✅ Order confirmed

        else Payment fails
            PAY -->> ORC : charged=false, reason=DECLINED
            Note over ORC: Begin compensation
            ORC ->> INV : releaseReservation(R99)
            INV -->> ORC : released
            ORC ->> ORD : updateStatus(CANCELLED)
            ORD -->> ORC : ❌ Order cancelled
        end

    else Items unavailable
        INV -->> ORC : reserved=false, reason=OUT_OF_STOCK
        ORC ->> ORD : updateStatus(FAILED)
        ORD -->> ORC : ❌ Order failed — stock
    end
\end{mermaid}

% =============================================================
\section{Gantt Charts}
% =============================================================

Gantt charts map tasks onto a timeline with dependencies,
milestones, and parallel work streams.  The \term{sunset} theme
gives warm, high-contrast bars ideal for printed schedules.

% ── 6. Full product launch plan ───────────────────────────────
\begin{mermaid}[
    type    = gantt,
    theme   = sunset,
    caption = {SaaS product launch — 16-week roadmap with milestones},
    label   = {fig:launch},
    width   = \linewidth
]
gantt
    title       SaaS Product Launch — Q3/Q4 2026
    dateFormat  YYYY-MM-DD
    excludes    weekends

    section Discovery
        Stakeholder interviews      :done,     s1, 2026-07-01, 10d
        Competitor analysis         :done,     s2, after s1,  8d
        Requirements workshop       :done,     s3, after s2,  5d
        Spec freeze                 :milestone,    after s3,  0d

    section Design
        Information architecture    :done,     d1, 2026-07-20, 8d
        Low-fi wireframes           :done,     d2, after d1,  6d
        High-fi UI design           :active,   d3, after d2,  12d
        Design system tokens        :active,   d4, 2026-08-01, 10d
        Usability testing           :          d5, after d3,  5d
        Design sign-off             :milestone,    after d5,  0d

    section Engineering — Backend
        Database schema             :done,     e1, 2026-07-15, 6d
        Auth service                :done,     e2, after e1,  8d
        Core API v1                 :active,   e3, after e2,  14d
        Payment integration         :          e4, after e3,  10d
        Admin dashboard API         :          e5, after e4,  8d

    section Engineering — Frontend
        Component library           :active,   f1, 2026-08-04, 10d
        Auth UI flows               :          f2, after f1,  6d
        Main product screens        :          f3, after f2,  12d
        Billing & settings UI       :          f4, after f3,  8d

    section QA & Launch
        Internal beta               :crit,     q1, 2026-09-22, 8d
        Beta feedback & fixes       :crit,     q2, after q1,  10d
        Performance testing         :crit,     q3, after q2,  5d
        Security audit              :crit,     q4, 2026-10-10, 7d
        Soft launch                 :milestone,    2026-10-20, 0d
        Marketing push              :          m1, 2026-10-20, 14d
        Public launch               :milestone,    2026-11-03, 0d
\end{mermaid}

% =============================================================
\section{Class Diagrams}
% =============================================================

Class diagrams express object-oriented domain models — classes,
attributes, methods, and their relationships.

% ── 7. Hospital management domain ────────────────────────────
\begin{mermaid}[
    type    = class,
    theme   = arctic,
    caption = {Hospital management system — domain class model},
    label   = {fig:hospital},
    width   = \linewidth
]
classDiagram
    direction TB

    class Person {
        <<abstract>>
        +int id
        +string firstName
        +string lastName
        +Date dateOfBirth
        +string email
        +fullName() string
        +age() int
    }

    class Patient {
        +string nhsNumber
        +BloodType bloodType
        +List~Allergy~ allergies
        +register()
        +getRecords() List~MedicalRecord~
    }

    class Doctor {
        +string licenseNumber
        +string specialisation
        +List~Qualification~ qualifications
        +isAvailable(Date) bool
        +prescribe(Patient, Drug) Prescription
    }

    class Nurse {
        +string ward
        +string shift
        +administerMedication(Patient, Prescription)
    }

    class Appointment {
        +int id
        +DateTime scheduledAt
        +string status
        +string notes
        +cancel()
        +reschedule(DateTime)
    }

    class MedicalRecord {
        +int id
        +Date date
        +string diagnosis
        +string treatment
        +List~Prescription~ prescriptions
    }

    class Prescription {
        +string drugName
        +float dosageMg
        +int frequencyPerDay
        +int durationDays
        +isExpired() bool
    }

    class Ward {
        +string name
        +int capacity
        +int currentOccupancy
        +hasAvailableBed() bool
        +admitPatient(Patient)
    }

    Person  <|-- Patient      : extends
    Person  <|-- Doctor       : extends
    Person  <|-- Nurse        : extends
    Patient "1" --> "0..*" Appointment   : books
    Doctor  "1" --> "0..*" Appointment   : attends
    Patient "1" *-- "0..*" MedicalRecord : has
    Doctor  "1" --> "0..*" Prescription  : writes
    MedicalRecord "1" *-- "0..*" Prescription : includes
    Nurse   "many" --> "1" Ward          : works in
    Patient "0..*" --> "0..1" Ward       : admitted to
\end{mermaid}

% =============================================================
\section{State Diagrams}
% =============================================================

State diagrams show how an entity moves through states in
response to events.  Nested states cleanly capture sub-processes.

% ── 8. Subscription billing lifecycle ────────────────────────
\begin{mermaid}[
    type    = state,
    theme   = forest,
    caption = {SaaS subscription billing state machine with nested states},
    label   = {fig:subscription},
    width   = 0.90\linewidth
]
stateDiagram-v2
    [*] --> Onboarding : User signs up

    state Onboarding {
        [*]             --> EmailVerification
        EmailVerification --> ProfileSetup : Email confirmed
        ProfileSetup    --> [*]            : Profile complete
    }

    Onboarding --> FreeTrial : Profile complete

    state FreeTrial {
        [*]          --> TrialActive
        TrialActive  --> TrialExpired : 14 days elapsed
        TrialActive  --> [*]         : User upgrades
        TrialExpired --> [*]
    }

    FreeTrial --> Active       : Subscribed (trial or direct)
    FreeTrial --> Churned      : Trial ended, no upgrade

    state Active {
        [*]        --> CurrentPeriod
        CurrentPeriod --> Renewing   : Period ends, payment OK
        Renewing   --> CurrentPeriod : Renewed
        CurrentPeriod --> PastDue    : Payment failed
    }

    Active  --> PastDue    : Payment failed
    PastDue --> Active     : Payment recovered (grace 7d)
    PastDue --> Suspended  : Grace period expired

    Suspended --> Active   : Payment recovered
    Suspended --> Churned  : 30 days suspended

    Active    --> Cancelled : User cancels
    Cancelled --> Active    : User re-subscribes (within 30d)
    Cancelled --> Churned   : 30 days passed

    Churned   --> [*]
\end{mermaid}

% =============================================================
\section{Entity-Relationship Diagrams}
% =============================================================

ER diagrams specify database schemas — entities, attributes,
primary and foreign keys, and cardinality.

% ── 9. Multi-tenant SaaS database ────────────────────────────
\begin{mermaid}[
    type    = er,
    theme   = midnight,
    caption = {Multi-tenant SaaS platform — full database schema},
    label   = {fig:schema},
    width   = \linewidth
]
erDiagram
    ORGANISATION {
        uuid   id           PK
        string name
        string slug
        string plan
        date   createdAt
        bool   isActive
    }
    USER {
        uuid   id           PK
        uuid   orgId        FK
        string email
        string hashedPwd
        string role
        date   lastLoginAt
    }
    PROJECT {
        uuid   id           PK
        uuid   orgId        FK
        uuid   ownerId      FK
        string name
        string status
        date   deadline
    }
    TASK {
        uuid   id           PK
        uuid   projectId    FK
        uuid   assigneeId   FK
        string title
        string priority
        string status
        date   dueDate
    }
    COMMENT {
        uuid   id           PK
        uuid   taskId       FK
        uuid   authorId     FK
        text   body
        date   createdAt
    }
    LABEL {
        uuid   id           PK
        uuid   orgId        FK
        string name
        string colourHex
    }
    ATTACHMENT {
        uuid   id           PK
        uuid   taskId       FK
        string fileName
        string s3Key
        int    sizeBytes
    }
    AUDIT_LOG {
        uuid   id           PK
        uuid   userId       FK
        string action
        json   payload
        date   occurredAt
    }

    ORGANISATION ||--o{ USER       : "has members"
    ORGANISATION ||--o{ PROJECT    : "owns"
    ORGANISATION ||--o{ LABEL      : "defines"
    USER         ||--o{ PROJECT    : "owns"
    PROJECT      ||--o{ TASK       : "contains"
    USER         ||--o{ TASK       : "assigned"
    TASK         ||--o{ COMMENT    : "has"
    USER         ||--o{ COMMENT    : "writes"
    TASK         }o--o{ LABEL      : "tagged"
    TASK         ||--o{ ATTACHMENT : "has"
    USER         ||--o{ AUDIT_LOG  : "generates"
\end{mermaid}

% =============================================================
\section{Timeline}
% =============================================================

% ── 10. History of programming languages ─────────────────────
\begin{mermaid}[
    type    = timeline,
    theme   = ocean,
    caption = {History of programming languages — key milestones},
    label   = {fig:lang-history},
    width   = \linewidth
]
timeline
    title  The History of Programming Languages
    section 1950s
        1957 : Fortran — first compiled language
             : COBOL design begins
        1958 : LISP — symbolic computation
             : ALGOL 58 — structured programming
    section 1960s–70s
        1964 : BASIC — beginners' language
        1972 : C — systems programming
             : Prolog — logic programming
        1975 : Scheme — minimalist LISP
    section 1980s
        1980 : Smalltalk 80 — objects everywhere
        1983 : C++ — C with classes
        1987 : Perl — text processing powerhouse
        1989 : Python conception
    section 1990s
        1991 : Python 0.9 released publicly
        1993 : Ruby created by Matz
        1995 : Java — write once, run anywhere
             : JavaScript — born in 10 days
             : PHP — web scripting
    section 2000s
        2001 : C# and .NET framework
        2003 : Scala — functional JVM language
        2007 : Clojure — LISP for the JVM
        2009 : Go — concurrency for cloud
             : Node.js — JS on the server
    section 2010s
        2011 : Kotlin — modern JVM language
        2012 : TypeScript — typed JavaScript
        2014 : Swift — Apple's modern language
             : Rust — memory-safe systems lang
    section 2020s
        2020 : Zig gaining traction
        2022 : Carbon — experimental C++ successor
        2024 : Mojo — AI/ML systems language
        2025 : AI-assisted language design era
\end{mermaid}

% =============================================================
\section{Mind Maps}
% =============================================================

% ── 11. Cloud architecture mindmap ────────────────────────────
\begin{mermaid}[
    type    = mindmap,
    theme   = forest,
    caption = {Cloud-native architecture concept map},
    label   = {fig:cloudmind},
    width   = \linewidth
]
mindmap
  root((Cloud Native Architecture))
    Compute
      Containers
        Docker
        Podman
      Orchestration
        Kubernetes
        Nomad
      Serverless
        AWS Lambda
        Cloud Run
        Azure Functions
    Networking
      Service Mesh
        Istio
        Linkerd
      API Gateway
        Kong
        AWS API GW
      CDN
        Cloudflare
        Fastly
    Data
      Relational
        PostgreSQL
        MySQL
        CockroachDB
      NoSQL
        MongoDB
        Cassandra
        DynamoDB
      Streaming
        Kafka
        Kinesis
        Pulsar
      Cache
        Redis
        Memcached
    Observability
      Metrics
        Prometheus
        Grafana
      Logs
        Loki
        ELK Stack
      Traces
        Jaeger
        Tempo
        Zipkin
    Security
      Identity
        Keycloak
        Auth0
        Okta
      Secrets
        Vault
        AWS Secrets Manager
      Policy
        OPA
        Kyverno
    CI/CD
      Build
        GitHub Actions
        GitLab CI
        Tekton
      GitOps
        ArgoCD
        Flux
      Artifact
        Harbor
        ECR
        GHCR
\end{mermaid}

% =============================================================
\section{Pie Charts}
% =============================================================

% ── 12. Dev time allocation ───────────────────────────────────
\begin{mermaid}[
    type    = pie,
    theme   = sunset,
    caption = {Developer time allocation across engineering activities},
    label   = {fig:devtime},
    width   = 0.75\linewidth
]
pie showData
    title Developer Time Allocation — Engineering Survey 2026
    "Feature development"     : 31.4
    "Bug fixing"              : 18.7
    "Code review"             : 14.2
    "Meetings and planning"   : 13.8
    "Documentation"           :  7.9
    "Testing"                 :  7.1
    "Refactoring / tech debt" :  4.5
    "On-call & incidents"     :  2.4
\end{mermaid}

% =============================================================
\section{Git Graphs}
% =============================================================

Git graphs visualise branching strategy, commit history,
merges, and release tags — ideal for documenting team workflows.

% ── 13. GitFlow branching strategy ───────────────────────────
\begin{mermaid}[
    type    = git,
    theme   = arctic,
    caption = {GitFlow branching strategy with feature branches and hotfixes},
    label   = {fig:gitflow},
    width   = \linewidth
]
gitGraph
    commit id: "Initial project scaffold"
    commit id: "Add CI/CD pipeline"
    commit id: "Configure linting & formatting"

    branch develop
    checkout develop
    commit id: "Set up dev environment"
    commit id: "Add database migrations"

    branch feature/user-auth
    checkout feature/user-auth
    commit id: "Add user model"
    commit id: "JWT middleware"
    commit id: "Login & register endpoints"
    commit id: "Password reset flow"

    checkout develop
    merge feature/user-auth id: "Merge: user auth"

    branch feature/dashboard
    checkout feature/dashboard
    commit id: "Dashboard layout"
    commit id: "Analytics widgets"
    commit id: "Export to CSV"

    branch feature/payments
    checkout feature/payments
    commit id: "Stripe SDK setup"
    commit id: "Checkout flow"
    commit id: "Webhooks handler"
    commit id: "Subscription model"

    checkout develop
    merge feature/payments id: "Merge: payments"
    merge feature/dashboard id: "Merge: dashboard"
    commit id: "Integration test suite"

    checkout main
    merge develop id: "Release v1.0.0" tag: "v1.0.0"

    branch hotfix/security-patch
    checkout hotfix/security-patch
    commit id: "Patch XSS vulnerability" type: HIGHLIGHT
    commit id: "Add CSP headers"

    checkout main
    merge hotfix/security-patch id: "Hotfix v1.0.1" tag: "v1.0.1"

    checkout develop
    merge hotfix/security-patch id: "Backport security fix"

    branch feature/mobile-api
    checkout feature/mobile-api
    commit id: "REST API v2 design"
    commit id: "Pagination & filtering"
    commit id: "Rate limiting"

    checkout develop
    merge feature/mobile-api id: "Merge: mobile API"
    commit id: "Load testing passed"
    commit id: "Docs updated"

    checkout main
    merge develop id: "Release v1.1.0" tag: "v1.1.0"
\end{mermaid}

% =============================================================
\section{User Journey}
% =============================================================

% ── 14. E-commerce purchase journey ──────────────────────────
\begin{mermaid}[
    type    = journey,
    theme   = ocean,
    caption = {Customer journey — discovering and purchasing a product online},
    label   = {fig:journey},
    width   = \linewidth
]
journey
    title Customer Purchase Journey — E-commerce Platform

    section Discovery
        See social media ad      : 3 : Customer
        Click ad to landing page : 4 : Customer
        Browse product catalogue : 4 : Customer, System
        Use search and filters   : 5 : Customer, System
        View product detail page : 5 : Customer

    section Evaluation
        Read product description : 4 : Customer
        Check reviews and ratings: 5 : Customer
        View photos and video    : 5 : Customer, System
        Compare with alternatives: 3 : Customer
        Check delivery estimate  : 4 : Customer, System

    section Purchase
        Add to cart              : 5 : Customer
        View cart                : 4 : Customer
        Enter shipping address   : 3 : Customer
        Choose delivery option   : 4 : Customer
        Enter payment details    : 2 : Customer
        Review and confirm order : 4 : Customer
        Payment processed        : 5 : Customer, Payment Service

    section Post-Purchase
        Receive order confirmation email : 5 : Customer, System
        Track shipment               : 4 : Customer, System
        Item delivered               : 5 : Customer, Courier
        Leave review                 : 4 : Customer
        Contact support (question)   : 2 : Customer, Support
\end{mermaid}

% =============================================================
\section{All Themes Side by Side}
% =============================================================

The same diagram rendered in each of the five built-in themes,
showing how \pkg{\textbackslash MermaidSetTheme\{\}} switches
the global default between sections.

% ── 15–19: theme showcase ─────────────────────────────────────

% ocean (already the global default)
\begin{mermaid}[
    type    = flowchart,
    theme   = ocean,
    caption = {Theme \texttt{ocean} — dark navy with electric-blue accents},
    label   = {fig:theme-ocean},
    width   = 0.80\linewidth
]
flowchart LR
    A[/"📥 Input"/] --> B{"Valid?"}
    B -->|✅ Yes| C[/"⚙️ Process"/] --> D[/"📤 Output"/]
    B -->|❌ No|  E[/"⚠️ Error"/]  --> F(["🔄 Retry"])
    F --> A
\end{mermaid}

\begin{mermaid}[
    type    = flowchart,
    theme   = forest,
    caption = {Theme \texttt{forest} — deep green, natural palette},
    label   = {fig:theme-forest},
    width   = 0.80\linewidth
]
flowchart LR
    A[/"📥 Input"/] --> B{"Valid?"}
    B -->|✅ Yes| C[/"⚙️ Process"/] --> D[/"📤 Output"/]
    B -->|❌ No|  E[/"⚠️ Error"/]  --> F(["🔄 Retry"])
    F --> A
\end{mermaid}

\begin{mermaid}[
    type    = flowchart,
    theme   = sunset,
    caption = {Theme \texttt{sunset} — warm charcoal with amber highlights},
    label   = {fig:theme-sunset},
    width   = 0.80\linewidth
]
flowchart LR
    A[/"📥 Input"/] --> B{"Valid?"}
    B -->|✅ Yes| C[/"⚙️ Process"/] --> D[/"📤 Output"/]
    B -->|❌ No|  E[/"⚠️ Error"/]  --> F(["🔄 Retry"])
    F --> A
\end{mermaid}

\begin{mermaid}[
    type    = flowchart,
    theme   = arctic,
    caption = {Theme \texttt{arctic} — clean white with slate-blue accents},
    label   = {fig:theme-arctic},
    width   = 0.80\linewidth
]
flowchart LR
    A[/"📥 Input"/] --> B{"Valid?"}
    B -->|✅ Yes| C[/"⚙️ Process"/] --> D[/"📤 Output"/]
    B -->|❌ No|  E[/"⚠️ Error"/]  --> F(["🔄 Retry"])
    F --> A
\end{mermaid}

\begin{mermaid}[
    type    = flowchart,
    theme   = midnight,
    caption = {Theme \texttt{midnight} — near-black with violet accents},
    label   = {fig:theme-midnight},
    width   = 0.80\linewidth
]
flowchart LR
    A[/"📥 Input"/] --> B{"Valid?"}
    B -->|✅ Yes| C[/"⚙️ Process"/] --> D[/"📤 Output"/]
    B -->|❌ No|  E[/"⚠️ Error"/]  --> F(["🔄 Retry"])
    F --> A
\end{mermaid}

% =============================================================
\section{Quick Reference}
% =============================================================

\subsection{Environment syntax}

\begin{verbatim}
\begin{mermaid}[
    type    = <diagram-type>,   % badge label
    theme   = <theme-name>,     % per-diagram colour theme
    caption = {Your caption},   % figure caption
    label   = {fig:key},        % \label key for \ref
    width   = 0.85\linewidth,   % listing/image width (default)
    pos     = H,                % float placement
    image   = {figures/d.pdf},  % if set → image mode
]
... mermaid source ...
\end{mermaid}
\end{verbatim}

\subsection{Available diagram types}

\begin{center}
\renewcommand{\arraystretch}{1.3}
\begin{tabular}{lll}
\toprule
\textbf{Option value} & \textbf{Badge} & \textbf{First keyword in source} \\
\midrule
\texttt{flowchart} & FLOW  & \texttt{flowchart TD} / \texttt{graph LR} \\
\texttt{sequence}  & SEQ   & \texttt{sequenceDiagram} \\
\texttt{gantt}     & GANTT & \texttt{gantt} \\
\texttt{class}     & CLASS & \texttt{classDiagram} \\
\texttt{state}     & STATE & \texttt{stateDiagram-v2} \\
\texttt{er}        & ER    & \texttt{erDiagram} \\
\texttt{pie}       & PIE   & \texttt{pie} \\
\texttt{timeline}  & TIME  & \texttt{timeline} \\
\texttt{mindmap}   & MIND  & \texttt{mindmap} \\
\texttt{git}       & GIT   & \texttt{gitGraph} \\
\texttt{journey}   & UX    & \texttt{journey} \\
\bottomrule
\end{tabular}
\end{center}

\subsection{Listing of all figures}

\listoffigures

\bigskip
\MermaidLiveNote

\end{document}
